\chapter{Comment MTO et DDMRP se complètent}\label{ch:mto-ddmrp}

DDMRP ne remplace pas MTO. Les deux mécanismes répondent à deux questions différentes.
DDMRP répond à la question : \emph{combien de matière faut-il protéger ou réapprovisionner ?}
MTO répond à la question : \emph{quelle commande doit recevoir et consommer cette matière ?}

\begin{table}[H]
  \centering
  \caption{Responsabilités respectives de MTO et de DDMRP dans le modèle.}\label{tab:mto-ddmrp}
  \begin{tabularx}{\textwidth}{@{}>{\bfseries}L{3.6cm}YY@{}}
    \toprule
    Fonction & MTO & DDMRP \\
    \midrule
    Déclencheur & commande client & position de flux net sous le seuil \\
    Besoin & calculé par commande, selon la nomenclature & demande qualifiée agrégée par matière \\
    Réservation & engagement de matière par commande, atomique & aucune : DDMRP n'engage pas de matière \\
    Open Supply & alloue une part d'une réception à la commande & compte la totalité des réceptions ouvertes \\
    Anticipation & aucune : réagit à une commande & oui : protège la matière par des zones \\
    Décision de réapprovisionnement & ouvre au plus une réception pour le manque non couvert & ouvre une réception vers le TOG \\
    Exécution Make & lance la production une fois la matière réservée & non concerné \\
    \bottomrule
  \end{tabularx}
\end{table}

Cette répartition évite deux erreurs. La première serait de laisser DDMRP décider de l'ordre de
production : une commande perdrait alors son lien avec la matière qu'elle consomme. La seconde
serait de laisser une commande ouvrir un approvisionnement sans regarder l'Open Supply : la même
quantité serait commandée deux fois. La coordination passe par la NFP et par l'allocation.

La demande d'une commande entre dans la demande qualifiée tant qu'elle n'est pas satisfaite
physiquement. Une réservation ne la retire donc pas de la demande qualifiée ; la consommation, elle,
la retire. Cette convention garantit qu'une même demande n'est soustraite qu'une fois. La
réservation reste un mécanisme d'intégrité physique : elle empêche d'engager deux fois la même
matière, mais elle ne modifie pas la définition DDMRP de la demande.

Une fois ces règles posées, le cas expérimental principal permet de voir l'ensemble en action
(\cref{ch:cas}).
